\documentclass{article}
\usepackage[UTF8]{ctex}
\usepackage{amsmath,amssymb,amsfonts}
\usepackage{enumitem}
\usepackage{geometry}
\geometry{a4paper, margin=1in}
% === 错题本排版优化（不参与题目标识） ===
\setlength{\parindent}{0pt}
\setlength{\parskip}{0.35em plus 0.10em minus 0.08em}
\setlength{\emergencystretch}{3em}
\setlist{itemsep=0.25em, topsep=0.40em, parsep=0pt}
\allowdisplaybreaks[2]
% === 错题本排版优化结束 ===
\begin{document}
\section*{中国海洋大学2024年高等代数真题}
\section*{116．中国海洋大学2024年高等代数试卷}

\subsection{题目1}
1．设 \(f_{1}(x), f_{2}(x), f_{3}(x)\) 是数域 \(P\) 上的多项式，证明：

\[
x^{3}+x^{2}+x+1 \mid x^{2} f_{1}\left(x^{4}\right)+x f_{2}\left(x^{4}\right)+f_{3}\left(x^{4}\right)
\]

当且仅当 \((x-1) \mid f_{i}(x), i=1,2,3\) ．

\subsection{题目2}
2．已知 \(\alpha_{1}, \alpha_{2}, \cdots, \alpha_{t}\) 是方程组 \(A X=0\) 的基础解系，证明：

\[
2 \alpha_{1}+\alpha_{2}, 2 \alpha_{2}+\alpha_{3}, \cdots, 2 \alpha_{t-1}+\alpha_{t}, 2 \alpha_{t}+\alpha_{1}
\]

也是方程组 \(A X=0\) 的基础解系。

\subsection{题目3}
3．已知 \(A\) 是 \(n\) 阶方阵，多项式 \(g(x)\) 满足 \(g(A)=O\) 。证明：对于任意非零多项式 \(f(\lambda)\) ，当\((f(\lambda), g(\lambda))=d(\lambda)\) 时，\(f(A)\) 与 \(d(A)\) 的秩相等．

\subsection{题目4}
4．已知

\[
W_{1}=\left\{\left.\left(\begin{array}{ccc}
a & 0 & d \\
c & 0 & b \\
0 & c & 0
\end{array}\right) \right\rvert\, a, b, c, d \in \mathbb{R}\right\}, W_{2}=\left\{\left.\left(\begin{array}{ccc}
x & 0 & z \\
0 & y & 0 \\
0 & 0 & 0
\end{array}\right) \right\rvert\, x, y, z \in \mathbb{R}\right\} .
\]

求 \(W_{1}+W_{2}\) 与 \(W_{1} \cap W_{2}\) ，并指出维数和基。

\subsection{题目5}
5．已知线性空间 \(V, W_{1}, W_{2}\) 满足 \(V=W_{1} \oplus W_{2}\) ，对任意的 \(X \in V\) ，记 \(X=X_{1}+X_{2}, X_{i} \in W_{i}(i=1,2)\)，定义线性变换 \(\sigma\) ，使得 \(\sigma(X)=X_{1}\) ，则称 \(\sigma\) 是 \(V\) 上的投影变换。证明：线性变换 \(\sigma\) 是 \(V\) 的投影变换当且仅当 \(\sigma\) 在任意基下的矩阵 \(A\) 满足 \(A^{2}=A\) ．

\subsection{题目6}
6．已知矩阵

\[
A=\left(\begin{array}{ccc}
0 & 1 & 0 \\
0 & 0 & 1 \\
-2 & 3 & -1
\end{array}\right) .
\]

\begin{itemize}
  \item[（1）] 若 \(A\) 是有理数域上的矩阵，判断 \(A\) 是否可以对角化，并说明理由．
  \item[（2）] 若 \(A\) 是复数域上的矩阵，判断 \(A\) 是否可以对角化，并说明理由．
\end{itemize}

\subsection{题目7}
7．设 \(\sigma, \tau\) 是欧氏空间 \(V\) 上的两个线性变换，满足 \((\sigma \alpha, \beta)=(\alpha, \tau \beta), \forall \alpha, \beta \in V\) 。证明：\(\sigma\) 的核是 \(\tau\) 的值域的正交补空间。

\subsection{题目8}
8．\(A\) 是 \(n\) 阶反对称阵，证明：存在非零列向量 \(\alpha\) ，使得 \(\left|A+\lambda \alpha \alpha^{\prime}\right|=|A|\) ，其中 \(\lambda\) 为任意常数。
\end{document}
